\section{Diccionario de ingeniería y revisión documental de migración}
La estimación distingue las herramientas reproducibles de la revisión humana de antecedentes. La Tabla~\ref{tab:mig-cargas} desarrolla 28 paquetes de ingeniería, incluidos los cuatro de diagnóstico ya descritos. La numeración final $c$ de los paquetes 1.8.2.$c$, 1.8.3.$c$ y 1.8.4.$c$ identifica, respectivamente, socios y armadores; contratos vigentes y documentos; seis años de pagos y saldos; embarcaciones e invernadas; los 14 expedientes de abandono; y tres años de matrícula escolar. Cada paquete produce un resultado propio de su conjunto.

\begingroup
\setlength{\tabcolsep}{2pt}
\begin{longtable}{L{\dimexpr.15\linewidth-2\tabcolsep\relax}L{\dimexpr.43\linewidth-2\tabcolsep\relax}R{\dimexpr.06\linewidth-2\tabcolsep\relax}R{\dimexpr.06\linewidth-2\tabcolsep\relax}R{\dimexpr.06\linewidth-2\tabcolsep\relax}R{\dimexpr.06\linewidth-2\tabcolsep\relax}R{\dimexpr.06\linewidth-2\tabcolsep\relax}R{\dimexpr.07\linewidth-2\tabcolsep\relax}R{\dimexpr.05\linewidth-2\tabcolsep\relax}}
\caption{Cargas iniciales de ingeniería de migración por paquete}\label{tab:mig-cargas}\\
\hline
 \EncabezadoTabla{Código} & \EncabezadoTabla{Resultado} & \EncabezadoTabla{A} & \EncabezadoTabla{D} & \EncabezadoTabla{S} & \EncabezadoTabla{Q} & \EncabezadoTabla{V} & \EncabezadoTabla{HH} & \EncabezadoTabla{Mes}\\\hline
\endfirsthead
\hline
 \EncabezadoTabla{Código} & \EncabezadoTabla{Resultado} & \EncabezadoTabla{A} & \EncabezadoTabla{D} & \EncabezadoTabla{S} & \EncabezadoTabla{Q} & \EncabezadoTabla{V} & \EncabezadoTabla{HH} & \EncabezadoTabla{Mes}\\\hline
\endhead
1.8.1.1 & Catálogo de fuentes y accesos & 16 & 8 & 0 & 0 & 0 & 24 & 1 \\\hline
1.8.1.2 & Extracción diagnóstica reproducible & 0 & 32 & 8 & 0 & 0 & 40 & 2\\\hline
1.8.1.3 & Informe de calidad del origen & 0 & 24 & 0 & 8 & 0 & 32 & 2 \\\hline
1.8.1.4 & Matriz de cobertura y decisiones & 8 & 8 & 0 & 8 & 0 & 24 & 3\\\hline
1.8.2.1 & Reglas: Socios y armadores & 8 & 16 & 0 & 0 & 0 & 24 & 3 \\\hline
1.8.2.2 & Reglas: Contratos vigentes y documentos & 8 & 16 & 0 & 0 & 0 & 24 & 3\\\hline
1.8.2.3 & Reglas: Seis años de pagos y saldos & 8 & 16 & 0 & 0 & 0 & 24 & 3 \\\hline
1.8.2.4 & Reglas: Embarcaciones e invernadas & 8 & 16 & 0 & 0 & 0 & 24 & 3\\\hline
1.8.2.5 & Reglas: 14 expedientes de abandono & 8 & 16 & 0 & 0 & 0 & 24 & 3 \\\hline
1.8.2.6 & Reglas: Tres años de matrícula escolar & 8 & 16 & 0 & 0 & 0 & 24 & 3\\\hline
1.8.3.1 & Saneamiento automatizado: Socios y armadores & 4 & 24 & 0 & 4 & 0 & 32 & 7 \\\hline
1.8.3.2 & Saneamiento automatizado: Contratos vigentes y documentos & 4 & 24 & 0 & 4 & 0 & 32 & 7\\\hline
1.8.3.3 & Saneamiento automatizado: Seis años de pagos y saldos & 4 & 24 & 0 & 4 & 0 & 32 & 7 \\\hline
1.8.3.4 & Saneamiento automatizado: Embarcaciones e invernadas & 4 & 24 & 0 & 4 & 0 & 32 & 7\\\hline
1.8.3.5 & Saneamiento automatizado: 14 expedientes de abandono & 4 & 24 & 0 & 4 & 0 & 32 & 7 \\\hline
1.8.3.6 & Saneamiento automatizado: Tres años de matrícula escolar & 4 & 24 & 0 & 4 & 0 & 32 & 7\\\hline
1.8.4.1 & Carga y conciliación: Socios y armadores & 0 & 24 & 0 & 8 & 8 & 40 & 8 \\\hline
1.8.4.2 & Carga y conciliación: Contratos vigentes y documentos & 0 & 24 & 0 & 8 & 8 & 40 & 8\\\hline
1.8.4.3 & Carga y conciliación: Seis años de pagos y saldos & 0 & 24 & 0 & 8 & 8 & 40 & 8 \\\hline
1.8.4.4 & Carga y conciliación: Embarcaciones e invernadas & 0 & 24 & 0 & 8 & 8 & 40 & 8\\\hline
1.8.4.5 & Carga y conciliación: 14 expedientes de abandono & 0 & 24 & 0 & 8 & 8 & 40 & 8 \\\hline
1.8.4.6 & Carga y conciliación: Tres años de matrícula escolar & 0 & 24 & 0 & 8 & 8 & 40 & 8\\\hline
1.8.5.1 & Repositorio e índice histórico & 0 & 8 & 4 & 4 & 24 & 40 & 9 \\\hline
1.8.5.2 & Consulta histórica con permisos & 0 & 8 & 4 & 4 & 24 & 40 & 9\\\hline
1.8.5.3 & Exportación y comprobación de integridad & 0 & 8 & 4 & 4 & 24 & 40 & 9 \\\hline
1.8.6.1 & Detección y carga idempotente de delta & 0 & 24 & 4 & 4 & 8 & 40 & 9\\\hline
1.8.6.2 & Herramientas de retorno previo & 0 & 24 & 4 & 4 & 8 & 40 & 9 \\\hline
1.8.6.3 & Herramientas de retorno posterior y diario & 0 & 24 & 4 & 4 & 8 & 40 & 9\\\hline
\textbf{Subtotal} & \textbf{Ingeniería de migración} & 96 & 552 & 32 & 112 & 144 & \textbf{936} & \\\hline
\end{longtable}
\FuenteTabla{Estimación ascendente de Synaptix. A: Análisis Funcional; D: Datos; S: Seguridad; Q: Calidad; V: Desarrollo. Horas-persona; meses relativos propuestos.}
\endgroup


\subsection{Correspondencias y transformaciones: 1.8.2.1--1.8.2.6}
Cada paquete entregará una especificación versionada de origen, destino, claves, relaciones y reglas del conjunto. Comprende 8 horas de Análisis Funcional para significado y validación, y 16 de Datos: 8 para correspondencias y 8 para reglas y ejemplos. La estimación de 24 horas por conjunto supone una fuente principal con hasta cinco formatos asociados y una ronda consolidada de decisión; formatos adicionales o reglas controvertidas se medirán al cerrar el diagnóstico.

El Líder de Datos responderá por el resultado. La recepción exigirá cobertura de los campos y vínculos obligatorios, casos de valores ausentes, destinos operativos o históricos y conformidad funcional de las decisiones que condicionen la carga. Las decisiones pendientes quedarán identificadas y bloquearán la transformación afectada. Se requieren el catálogo, el informe de calidad, el modelo de SD5 y los validadores del cliente. El hito será la recepción de la especificación del conjunto, anterior a su saneamiento y construcción de carga. Referencias: SD5, secciones 5.1 y 5.3; RNF-MIG-01; BT RT-05.11.

\subsection{Saneamiento automatizado: 1.8.3.1--1.8.3.6}
Cada paquete entregará reglas ejecutables y una copia saneada con registro de cambios, conservando el origen. Datos dedicará 16 horas a implementar reglas y 8 a ejecutar y documentar su resultado; Análisis Funcional dedicará 4 a decisiones y Calidad 4 a comprobar reproducibilidad. Las 32 horas por conjunto suponen hasta cinco clases de defecto tratables mediante reglas; no incluyen lectura manual exhaustiva, OCR ni recuperación de archivos dañados.

El Líder de Datos responderá por el resultado. La aceptación comprobará que cada corrección identifica regla, versión y registros afectados; los defectos no corregibles conservarán decisión y responsable. No se inventarán pagos ni autorizaciones. Se requieren copia protegida y reglas recibidas en 1.8.2. Su hito habilitará la prueba del cargador correspondiente. Referencias: SD5, secciones 5.2.4 y 5.3.2; BT RT-05.12.

\subsection{Herramientas de carga y conciliación: 1.8.4.1--1.8.4.6}
Cada paquete entregará un cargador y sus controles de lote para el conjunto correspondiente. Datos dedicará 16 horas a transformación y carga y 8 a manifiestos y conciliación; Desarrollo dedicará 8 a automatización y manejo de errores; Calidad dedicará 8 a pruebas locales de reejecución y controles. El esfuerzo será 40 horas por conjunto. Las herramientas utilizarán el diagnóstico y las interfaces especificadas; la estimación supone una ruta de carga principal por conjunto, sin incluir extracción por ingeniería inversa.

El Líder de Datos responderá por el resultado. La aceptación verificará correspondencias, recuentos, sumas monetarias separadas por moneda, integridad y reejecución sin efectos duplicados. La prueba local de un cargador no será un ensayo completo de migración. Se requieren reglas recibidas, copia saneada y esquema destino. El hito será la recepción del cargador para su integración en 1.10.1. Referencias: SD5, secciones 5.3.1--5.3.2; BT RT-05.13--RT-05.14.

\subsection{Archivo histórico: 1.8.5.1--1.8.5.3}
Los paquetes entregarán, respectivamente, repositorio e índice; consulta con permisos; y exportación con comprobación de integridad. Cada uno se estima en 40 horas: Desarrollo 24, Datos 8, Seguridad 4 y Calidad 4. Desarrollo distribuirá su trabajo en 16 horas de implementación y 8 de pruebas y documentación. La reutilización de almacenamiento, identidad y componentes de consulta de SD4 es un supuesto de estimación; su construcción base pertenece a otras ramas y no se suma nuevamente aquí.

El Líder de Datos responderá por los tres resultados, con ejecución de Desarrollo. La recepción del repositorio comprobará índice, relaciones y categorías de retención; la consulta comprobará perfiles, búsqueda y acceso a antecedentes; la exportación comprobará formatos abiertos, manifiestos e integridad. La consulta será independiente del ejecutable de 2014. La exportación de los 14 expedientes incluirá los controles PDF/A de CA-22; no se supondrá que la descarga de un archivo equivale a esa validación. Se requieren contratos de almacenamiento e identidad, y la clasificación de 1.8.2. Los hitos serán recepción del repositorio, de la consulta y de la exportación, anteriores a su ensayo integral. Referencias: SD5, secciones 5.2.5 y 5.3; BT RT-05.15; SD3 CA-22.

\subsection{Delta y retorno: 1.8.6.1--1.8.6.3}
Los paquetes entregarán detección y carga de altas, cambios y bajas; herramientas de retorno previo a nueva escritura; y herramientas de retorno posterior con diario conciliado. Cada paquete comprende 24 horas de Datos ---16 de implementación y 8 de pruebas locales---, 8 de Desarrollo para integración, 4 de Seguridad para permisos y protección y 4 de Calidad para revisión del resultado: 40 horas en total.

El Líder de Datos responderá por las herramientas. La aceptación del delta comprobará cambios sobre copias sucesivas sin presumir captura automática disponible. El retorno previo preservará la autoridad anterior y el diario; el posterior identificará hechos nuevos, respuestas externas pendientes y efectos que el legado no puede representar. Las herramientas no descartarán hechos para recuperar una copia anterior. Se requieren copia recuperable, correspondencias y contratos administrativos; la estimación supone que los cargadores ya están recibidos. La recepción habilitará los ensayos cronometrados de 1.10.1; las 40 horas no demuestran los umbrales de 90 o 120 minutos de SD5. Referencias: SD5, sección 5.3.3; BT RT-05.11 y RT-20.02.

\subsection{Revisión completa de expedientes y autorizaciones}
La revisión documental tendrá paquetes adicionales de tamaño controlable. Para cada expediente de abandono $e$, se define \texttt{1.8.3.7.e.b}, donde $b$ identifica un lote de hasta 24 piezas inventariadas. Cada lote se estima en 8 horas: Datos 2 para índice y vínculos, Análisis Funcional 4 para legibilidad y significado, y Calidad 2 para comprobación y evidencia. Estas tasas suponen piezas legibles y disponibles; digitalización, recuperación y aportación de antecedentes ausentes se tratarán por separado sin omitirlos del universo de recepción.

Sea $N_e$ la cantidad inventariada de piezas del expediente $e$. El número de lotes será $K=\sum_{e=1}^{14}\max(1,\lceil N_e/24\rceil)$. El lote de un expediente sin piezas entregará la constatación de ausencia y sus bloqueos; no acreditará un expediente completo. El esfuerzo propuesto será $8K$ horas. El Líder de Datos responderá por el registro; la aprobación del contenido probatorio corresponderá a Administración y la instancia competente del cliente. La recepción completa exigirá las piezas de los 14 expedientes con destino explicado y los controles de CA-22.

Para las autorizaciones escolares se define \texttt{1.8.3.8.b}, por lotes de 32--256 autorizaciones; un remanente inferior a 32 se incorporará a otro lote compatible. Se propone 0,25 horas por autorización: 0,15 de Análisis Funcional, 0,05 de Datos y 0,05 de Calidad. Con $N_A$ autorizaciones exigibles al corte, el esfuerzo será $0{,}25N_A$ horas. Si el universo completo es menor que 32, se formulará un paquete conjunto de revisión de al menos 8 horas que incluya sus decisiones y documentación, sin declarar una población ficticia.

El Líder de Datos responderá por los vínculos y el registro; Escuela validará vigencia y facultades de quienes autorizan. La aceptación verificará el 100\,\% de las autorizaciones que se utilizarán al corte y distinguirá antecedentes históricos de permisos vigentes. El inventario de tres años de matrícula no equivale a $N_A$. La revisión se reservará antes de cada activación escolar y se actualizará si cambian participantes o autorizaciones durante la marcha blanca; esa actualización pertenecerá a los paquetes de implantación. Referencia: SD5, secciones 5.3.1--5.3.3.

\subsection{Perímetro de la estimación}
Los 28 paquetes de ingeniería suman 936 horas-persona: Análisis Funcional 96, Datos 552, Seguridad 32, Calidad 112 y Desarrollo 144. El subtotal incluye las 120 horas de diagnóstico y sus pruebas locales. Con autorizaciones suficientes para aplicar los lotes indicados, el esfuerzo de esta parte de 1.8 será $936+8K+0{,}25N_A$; el caso de un universo escolar menor que 32 conservará el mínimo de paquete descrito. No se asigna todavía un valor observado a $K$ ni a $N_A$.

El subtotal excluye insumos, licencias no compartidas, respuesta a defectos fuera de los supuestos, los dos ensayos completos, los cortes de producción, marcha blanca, estabilización y operación. Los resultados del diagnóstico determinarán el número real de formatos, reglas y lotes y permitirán actualizar la estimación antes de cerrar la línea base. La oficina de proyecto registrará su impacto en costo, recursos y calendario sin modificar unilateralmente los hitos contractuales.
